\documentclass[12pt,a4paper]{article}
\usepackage[utf8]{inputenc}
\usepackage[indonesian]{babel}
\usepackage{amsmath}
\usepackage{amsfonts}
\usepackage{amssymb}
\usepackage{geometry}

\geometry{margin=2.5cm}

\title{Penyelesaian Titik Berat Benda Putar}
\author{Ahmad Fakhrul Bawani}
\date{}

\begin{document}

\maketitle

\section{Soal}
Tentukan titik berat benda yang dibatasi
\begin{equation}
  y = 4 - x^2
\end{equation}
diputar di sumbu-x di kuadran 1.

\section{Penyelesaian}

\subsection*{2.1 Menentukan Batas Integrasi dan Luas Daerah ($A$)}
Batas untuk $x$ adalah dari $x = 0$ hingga titik potong dengan sumbu-X ($y = 0$):
\[ 4 - x^2 = 0 \implies x^2 = 4 \implies x = 2 \]
Jadi, batas integrasi adalah $x \in [0, 2]$.

Luas daerah $A$ dihitung dengan:
\[ A = \int_{0}^{2} y \, dx = \int_{0}^{2} (4 - x^2) \, dx \]
\[ A = \left[ 4x - \frac{1}{3}x^3 \right]_{0}^{2} = \left( 4(2) - \frac{8}{3} \right) - 0 = 8 - \frac{8}{3} = \frac{16}{3} \]

\subsection*{2.2 Menghitung Koordinat $\bar{x}$}
Rumus untuk momen pertama terhadap sumbu-Y ($M_y$) dan koordinat $\bar{x}$ adalah:
\[ M_y = \int_{0}^{2} x \cdot y \, dx = \int_{0}^{2} x(4 - x^2) \, dx = \int_{0}^{2} (4x - x^3) \, dx \]
\[ M_y = \left[ 2x^2 - \frac{1}{4}x^4 \right]_{0}^{2} = \left( 2(2)^2 - \frac{1}{4}(2)^4 \right) - 0 = 8 - 4 = 4 \]

Maka, nilai $\bar{x}$ adalah:
\[ \bar{x} = \frac{M_y}{A} = \frac{4}{\frac{16}{3}} = 4 \times \frac{3}{16} = \frac{3}{4} = 0.75 \]

\subsection*{2.3 Menghitung Koordinat $\bar{y}$}
Rumus untuk momen pertama terhadap sumbu-X ($M_x$) dan koordinat $\bar{y}$ adalah:
\[ M_x = \frac{1}{2} \int_{0}^{2} y^2 \, dx = \frac{1}{2} \int_{0}^{2} (4 - x^2)^2 \, dx \]
\[ M_x = \frac{1}{2} \int_{0}^{2} (16 - 8x^2 + x^4) \, dx \]
\[ M_x = \frac{1}{2} \left[ 16x - \frac{8}{3}x^3 + \frac{1}{5}x^5 \right]_{0}^{2} \]
\[ M_x = \frac{1}{2} \left( 16(2) - \frac{8}{3}(2)^3 + \frac{1}{5}(2)^5 \right) = \frac{1}{2} \left( 32 - \frac{64}{3} + \frac{32}{5} \right) \]
\[ M_x = \frac{1}{2} \left( \frac{480 - 320 + 96}{15} \right) = \frac{1}{2} \left( \frac{256}{15} \right) = \frac{128}{15} \]

Maka, nilai $\bar{y}$ adalah:
\[ \bar{y} = \frac{M_x}{A} = \frac{\frac{128}{15}}{\frac{16}{3}} = \frac{128}{15} \times \frac{3}{16} = \frac{128}{16} \times \frac{3}{15} = 8 \times \frac{1}{5} = \frac{8}{5} = 1.6 \]

\subsection*{Kesimpulan Bukti Nilai $\bar{x}$ dan $\bar{y}$}
Dari hasil perhitungan di atas, kita mendapatkan:
\[ \bar{x} = \frac{3}{4} = 0.75 \]
\[ \bar{y} = \frac{8}{5} = 1.6 \]
Dengan demikian, terbukti secara matematis bahwa untuk daerah yang dibatasi oleh $y = 4 - x^2$ di kuadran 1:
\[ \bar{y} \neq \bar{x} \quad \left(1.6 \neq 0.75\right) \]

\end{document}